\SourceLeaf{176}\par\Needspace{6\baselineskip}\medskip
\begin{center}\hypertarget{native-176-article129}{}\pdfbookmark[2]{129.}{article-129}129.\end{center}
Aus den fünf bisher erläuterten Verfahren ergeben sich unmittelbar fünf weitere. Sie
unterscheiden sich von jenen allein dadurch, dass man für $x$ und $y$ die Bahnneigung und die Länge
des aufsteigenden Knotens anstelle der Abstände von der Erde einsetzt. Die neuen Verfahren laufen
folgendermaßen ab.

I. Aus $x$ und $y$ und den beiden äußeren geozentrischen Orten bestimmt man nach den Artikeln 74 und
75 die heliozentrischen Längen in der Bahn und die Radiusvektoren. Aus diesen und den entsprechenden
Zeiten bestimmt man alle übrigen Elemente. Aus diesen Elementen berechnet man schließlich den
geozentrischen Ort zum Zeitpunkt der mittleren Beobachtung. Die Differenzen seiner Länge und Breite
gegenüber Länge und Breite des beobachteten Ortes ergeben $X$ und $Y$.

Bei den vier übrigen Verfahren berechnet man zunächst aus der Lage der Bahnebene und den
geozentrischen Orten alle drei heliozentrischen Längen in der Bahn und die entsprechenden
Radiusvektoren. Danach geht man folgendermaßen vor.

II. Man bestimmt die übrigen Elemente allein aus den beiden äußeren Orten und den entsprechenden
Zeiten. Mithilfe dieser Elemente berechnet man für den Zeitpunkt der mittleren Beobachtung die Länge
in der Bahn und den Radiusvektor. Die Differenzen dieser Größen gegenüber den zuvor gefundenen, also
aus dem geozentrischen Ort abgeleiteten Werten ergeben $X,Y$.

III. Oder man leitet die übrigen Abmessungen der Bahn aus allen drei heliozentrischen Orten ab
(Artikel 82 und 83). In diese Rechnung gehen die Zeiten nicht ein. Anschließend ermittelt man,
welche Zeitintervalle auf der so gefundenen Bahn zwischen der ersten und der zweiten sowie zwischen
der zweiten und der dritten Beobachtung hätten verstreichen müssen. Die Differenzen dieser
Intervalle gegenüber den wahren Werten ergeben $X,Y$.

IV. Man berechnet die übrigen Elemente auf zweierlei Weise, zunächst aus der Kombination des ersten
mit dem zweiten Ort, dann aus der Kombination des zweiten mit dem dritten, jeweils unter Verwendung
der zugehörigen Zeitintervalle. Vergleicht man die beiden Systeme von Elementen miteinander, so kann
man zwei beliebige der Differenzen für $X,Y$ nehmen.

V. Schließlich kann man dieselbe doppelte Rechnung nur bis zu den Werten der in Artikel 91 mit $y$
bezeichneten Größe führen und für $X,Y$ die im vorigen Artikel angegebenen Ausdrücke verwenden.

Damit man die letzten vier Verfahren sicher anwenden kann, darf die Erde bei keiner der drei
Beobachtungen den Knoten der Bahn zu nahe sein. \SourceLeaf{177}Das erste Verfahren verlangt diese Bedingung nur
für die beiden äußeren Beobachtungen. Genauer gesagt darf höchstens einer der drei Orte der Erde in
der Nähe der Knoten liegen, da man den mittleren Ort anstelle eines der beiden äußeren verwenden
kann.
